\section{洪武至永乐的补充编年节点（一）}
这一组条目把此前尚未设导航入口的年度材料接回本章脉络。它们不是把同年材料机械等同为因果，而是在可核查的时间位置上，分别保留行动、行政记录与社会经验之间的距离。

\subsection{1368年前后的材料线索}

\eventanchor{EM-1368-DADU}{明·1368（洪武元年）八月}
\paragraph{1368（洪武元年）八月：徐达军进入大都并改称北平} 这一节点应放回本章已经展开的权力、财政、地方社会与边地关系中理解。账本把它出现的条件概括为：明军自山东、河南北上且元顺帝已北撤；大都守备难获外援 随后可追踪的变化是：元廷失去旧都和华北行政中心；明须经营北平并面对草原残余政权 可由《太祖实录》洪武元年八月条；《明史·太祖本纪》；《高丽史·恭愍王世家》二十一年条 互相核对；不过，日期可靠；明廷军报不独证城内损失和居民去留 因而这里标示的是材料能够支持的行动与制度位置，而不是对动机、地方执行或普通人经验的无保留复原。

\eventanchor{EM-1368-FOUNDING}{明·1368（洪武元年）正月}
\paragraph{1368（洪武元年）正月：朱元璋在应天即皇帝位并建国号明} 这一节点应放回本章已经展开的权力、财政、地方社会与边地关系中理解。账本把它出现的条件概括为：红巾军政权分裂而元廷南方控制崩解；朱元璋已控制江南粮赋与应天军政中心 随后可追踪的变化是：明廷取得颁诏和任官名义；北伐与地方接收须转为连续的行政和军事行动 可由《太祖实录》洪武元年正月条（即位诏）；《明史·太祖本纪》；南京明故宫遗址考古发掘报告 互相核对；不过，日期可靠；诏告可证建国程序而不能替代各地归附过程 因而这里标示的是材料能够支持的行动与制度位置，而不是对动机、地方执行或普通人经验的无保留复原。

\eventanchor{EM-1368-GUOZIJIAN}{明·1368（洪武元年）九月}
\paragraph{1368（洪武元年）九月：在南京设国子学以培训中央官员子弟} 这一节点应放回本章已经展开的权力、财政、地方社会与边地关系中理解。账本把它出现的条件概括为：新朝须为六部和地方官体系补充识字行政人才并重建礼制教育 随后可追踪的变化是：南京国子监成为明初官学枢纽；科举与荐举的关系仍待后续制度调整 可由《太祖实录》洪武元年九月条；《明会典·国子监》；《明史·选举志》 互相核对；不过，日期较可靠；学校建置不等于生员规模或教学成效 因而这里标示的是材料能够支持的行动与制度位置，而不是对动机、地方执行或普通人经验的无保留复原。

\eventanchor{EM-1369-SHANGDU}{明·1369（洪武二年）五月至八月}
\paragraph{1369（洪武二年）五月至八月：明军北进上都、应昌一线；元顺帝在北遁途中去世} 这一节点应放回本章已经展开的权力、财政、地方社会与边地关系中理解。账本把它出现的条件概括为：大都失守后北元仍试图依托上都、应昌和草原骑兵维持王朝中心 随后可追踪的变化是：明军未能消灭草原政治力量；北边战争由夺城转为长期机动与互市问题 可由《太祖实录》洪武二年五月、八月条；《明史·北虏传》；《高丽史》北元消息条 互相核对；不过，日期与战事序列较可靠；元顺帝死地和撤退细节有异说 因而这里标示的是材料能够支持的行动与制度位置，而不是对动机、地方执行或普通人经验的无保留复原。

\eventanchor{EM-1370-PRINCELY-ENFEOFFMENT}{明·1370（洪武三年）四月}
\paragraph{1370（洪武三年）四月：太祖大封诸子为王并配置藩国} 这一节点应放回本章已经展开的权力、财政、地方社会与边地关系中理解。账本把它出现的条件概括为：建国后皇室需要可继承的资源分配与边地防卫网络 随后可追踪的变化是：诸王成为地方军政力量；中央与藩王的权限张力为建文时期冲突埋下制度条件 可由《太祖实录》洪武三年四月条；《明史·诸王传》；《明会典·宗藩》 互相核对；不过，日期可靠；封册可证名义授封而不能直接证明各王实际兵权 因而这里标示的是材料能够支持的行动与制度位置，而不是对动机、地方执行或普通人经验的无保留复原。

\eventanchor{EM-1370-SHENER-BATTLE}{明·1370（洪武三年）六月}
\paragraph{1370（洪武三年）六月：徐达在沈儿峪击败扩廓帖木儿部} 这一节点应放回本章已经展开的权力、财政、地方社会与边地关系中理解。账本把它出现的条件概括为：明廷欲趁北元重组未定追击其主力；徐达军依托山西、陕西粮道北进 随后可追踪的变化是：扩廓部受挫但未被歼灭；明廷随后仍需以多路远征争夺草原主动权 可由《太祖实录》洪武三年六月条；《明史·徐达传》与《北虏传》；17世纪《蒙古源流》北元汗统叙事 互相核对；不过，战役发生可靠；双方兵数和斩获主要来自明方战报 因而这里标示的是材料能够支持的行动与制度位置，而不是对动机、地方执行或普通人经验的无保留复原。

\eventanchor{EM-1371-HAIJIN}{明·1371（洪武四年）十二月}
\paragraph{1371（洪武四年）十二月：朝廷申禁滨海居民私自出海并限制民间海外贸易} 这一节点应放回本章已经展开的权力、财政、地方社会与边地关系中理解。账本把它出现的条件概括为：倭患警报、元末海上武装遗绪与新朝对沿海人口流动的疑虑交织 随后可追踪的变化是：朝贡贸易被置于官府控制之下；私航、走私和沿海防务成为持续的地方—中央博弈 可由《太祖实录》洪武四年十二月条；《明史·食货志》；《筹海图编》所录早明海禁材料 互相核对；不过，诏令日期可靠；禁令不等于沿海航行和交易立即停止 因而这里标示的是材料能够支持的行动与制度位置，而不是对动机、地方执行或普通人经验的无保留复原。

\eventanchor{EM-1371-MINGXIA}{明·1371（洪武四年）}
\paragraph{1371（洪武四年）：傅友德等沿长江入蜀，夏主明昇降明} 这一节点应放回本章已经展开的权力、财政、地方社会与边地关系中理解。账本把它出现的条件概括为：明夏控制四川但与中原补给和政治盟友隔绝；明军掌握长江中下游军运条件 随后可追踪的变化是：四川纳入明朝行政与赋役框架；山地边区和土司关系仍须逐步处理 可由《太祖实录》洪武四年条；《明史·明玉珍传》；《四川通志》明夏遗迹条 互相核对；不过，年代可靠；行军路线和地方响应不能只据凯旋叙事复原 因而这里标示的是材料能够支持的行动与制度位置，而不是对动机、地方执行或普通人经验的无保留复原。

\eventanchor{ML-1371-003}{明·1371（洪武四年）八月3日}
\paragraph{1371（洪武四年）八月3日：明盛将四川投降明朝} 这一节点应放回本章已经展开的权力、财政、地方社会与边地关系中理解。账本把它出现的条件概括为：前期边防、地方武装或跨区域资源调动的压力推动了该行动。 随后可追踪的变化是：它改变了后续调兵、补给或地方秩序的条件；战果和伤亡须分开核对。 可由《太祖实录》洪武四年条；《明史·兵志》及相关列传；边镇、地方志或对方材料 互相核对；不过，译写候选以编年材料定位；实录记载反映朝廷选择，崇祯期须另以奏疏、地方及敌对政权材料交叉。 因而这里标示的是材料能够支持的行动与制度位置，而不是对动机、地方执行或普通人经验的无保留复原。

\eventanchor{ML-1371-004}{明·1371（洪武四年）}
\paragraph{1371（洪武四年）：国子监注册学生达3,728人} 这一节点应放回本章已经展开的权力、财政、地方社会与边地关系中理解。账本把它出现的条件概括为：继承、官僚协调或皇权运作的压力使问题进入朝廷程序。 随后可追踪的变化是：它影响后续人事与政策议程；动机和派系标签不能由单一叙事裁定。 可由《太祖实录》洪武四年条；《明史·本纪》及相关列传；诏令、奏疏或会典版本 互相核对；不过，译写候选以编年材料定位；实录记载反映朝廷选择，崇祯期须另以奏疏、地方及敌对政权材料交叉。 因而这里标示的是材料能够支持的行动与制度位置，而不是对动机、地方执行或普通人经验的无保留复原。

\eventanchor{EM-1372-LINGBEI-EXPEDITION}{明·1372（洪武五年）}
\paragraph{1372（洪武五年）：三路北伐受挫，李文忠、徐达等部未能取得预期战果} 这一节点应放回本章已经展开的权力、财政、地方社会与边地关系中理解。账本把它出现的条件概括为：明廷试图以多路大军压迫北元；行军距离、草原水草和部队协同限制远征 随后可追踪的变化是：太祖转向更审慎的北边经营；卫所、边墙和藩王防线的重要性上升 可由《太祖实录》洪武五年条；《明史·李文忠传》与《北虏传》；17世纪《蒙古源流》北元汗统叙事 互相核对；不过，战役大势可靠；败因、损失和各路会师情况在记载中不一 因而这里标示的是材料能够支持的行动与制度位置，而不是对动机、地方执行或普通人经验的无保留复原。

\eventanchor{EM-1372-RYUKYU}{明·1372（洪武五年）}
\paragraph{1372（洪武五年）：明使杨载出使琉球并开启朝贡往来} 这一节点应放回本章已经展开的权力、财政、地方社会与边地关系中理解。账本把它出现的条件概括为：海禁后明廷仍需以官方使节组织可控的海上外交；琉球各势力寻求外部承认 随后可追踪的变化是：形成持续的册封—贸易通道；朝贡名义与实际航运、商人网络不能混同 可由《太祖实录》洪武五年条；《明史·琉球传》；《历代宝案》早期入贡文书 互相核对；不过，明使出航与琉球回应有双边记录；三山内部政治不应视作单一国家 因而这里标示的是材料能够支持的行动与制度位置，而不是对动机、地方执行或普通人经验的无保留复原。

\eventanchor{ML-1372-001}{明·1372（洪武五年）四月}
\paragraph{1372（洪武五年）四月：明军在土拉河击败可克帖木儿} 这一节点应放回本章已经展开的权力、财政、地方社会与边地关系中理解。账本把它出现的条件概括为：财政征解、交通工程或货币支付的压力使相关措施进入政务。 随后可追踪的变化是：其法定安排不等于地方执行，后续须以地域材料追踪负担与成效。 可由《太祖实录》洪武五年条；《明会典》相关条目；地方志、黄册/契约/河工材料 互相核对；不过，译写候选以编年材料定位；实录记载反映朝廷选择，崇祯期须另以奏疏、地方及敌对政权材料交叉。 因而这里标示的是材料能够支持的行动与制度位置，而不是对动机、地方执行或普通人经验的无保留复原。

\eventanchor{ML-1372-002}{明·1372（洪武五年）}
\paragraph{1372（洪武五年）：明军在喀喇昆仑被击溃} 这一节点应放回本章已经展开的权力、财政、地方社会与边地关系中理解。账本把它出现的条件概括为：前期边防、地方武装或跨区域资源调动的压力推动了该行动。 随后可追踪的变化是：它改变了后续调兵、补给或地方秩序的条件；战果和伤亡须分开核对。 可由《太祖实录》洪武五年条；《明史·兵志》及相关列传；边镇、地方志或对方材料 互相核对；不过，译写候选以编年材料定位；实录记载反映朝廷选择，崇祯期须另以奏疏、地方及敌对政权材料交叉。 因而这里标示的是材料能够支持的行动与制度位置，而不是对动机、地方执行或普通人经验的无保留复原。

\eventanchor{ML-1372-003}{明·1372（洪武五年）}
\paragraph{1372（洪武五年）：明军攻克永昌，攻克居延} 这一节点应放回本章已经展开的权力、财政、地方社会与边地关系中理解。账本把它出现的条件概括为：前期边防、地方武装或跨区域资源调动的压力推动了该行动。 随后可追踪的变化是：它改变了后续调兵、补给或地方秩序的条件；战果和伤亡须分开核对。 可由《太祖实录》洪武五年条；《明史·兵志》及相关列传；边镇、地方志或对方材料 互相核对；不过，译写候选以编年材料定位；实录记载反映朝廷选择，崇祯期须另以奏疏、地方及敌对政权材料交叉。 因而这里标示的是材料能够支持的行动与制度位置，而不是对动机、地方执行或普通人经验的无保留复原。

\eventanchor{ML-1372-004}{明·1372（洪武五年）}
\paragraph{1372（洪武五年）：国子监注册学生突破1万人} 这一节点应放回本章已经展开的权力、财政、地方社会与边地关系中理解。账本把它出现的条件概括为：继承、官僚协调或皇权运作的压力使问题进入朝廷程序。 随后可追踪的变化是：它影响后续人事与政策议程；动机和派系标签不能由单一叙事裁定。 可由《太祖实录》洪武五年条；《明史·本纪》及相关列传；诏令、奏疏或会典版本 互相核对；不过，译写候选以编年材料定位；实录记载反映朝廷选择，崇祯期须另以奏疏、地方及敌对政权材料交叉。 因而这里标示的是材料能够支持的行动与制度位置，而不是对动机、地方执行或普通人经验的无保留复原。

\eventanchor{ML-1372-005}{明·1372（洪武五年）}
\paragraph{1372（洪武五年）：明代出现了专门为海军使用的大炮。} 这一节点应放回本章已经展开的权力、财政、地方社会与边地关系中理解。账本把它出现的条件概括为：制度、教育或知识传播的需要推动了相关行动或文本形成。 随后可追踪的变化是：它提供制度和传播史的锚点，不能据此推断社会整体接受。 可由《太祖实录》洪武五年条；《明史·选举志》或《艺文志》；文集、碑刻或书目材料 互相核对；不过，译写候选以编年材料定位；实录记载反映朝廷选择，崇祯期须另以奏疏、地方及敌对政权材料交叉。 因而这里标示的是材料能够支持的行动与制度位置，而不是对动机、地方执行或普通人经验的无保留复原。

\eventanchor{EM-1373-EXAMINATION-SUSPENDED}{明·1373（洪武六年）}
\paragraph{1373（洪武六年）：太祖停罢科举，改重荐举与学校考核} 这一节点应放回本章已经展开的权力、财政、地方社会与边地关系中理解。账本把它出现的条件概括为：首次科举所取士人与新朝行政需求之间发生张力；皇帝偏好直接控制人才来源 随后可追踪的变化是：荐举和学校渠道暂时扩张；科举何时、如何恢复成为中央选官制度的后续问题 可由《太祖实录》洪武六年条；《明史·选举志》；《明会典·礼部》科举沿革 互相核对；不过，政策日期较可靠；官方对取士质量的评语带有皇帝立场 因而这里标示的是材料能够支持的行动与制度位置，而不是对动机、地方执行或普通人经验的无保留复原。

\eventanchor{ML-1373-001}{明·1373（洪武六年）三月}
\paragraph{1373（洪武六年）三月：洪武皇帝暂停科举考试} 这一节点应放回本章已经展开的权力、财政、地方社会与边地关系中理解。账本把它出现的条件概括为：制度、教育或知识传播的需要推动了相关行动或文本形成。 随后可追踪的变化是：它提供制度和传播史的锚点，不能据此推断社会整体接受。 可由《太祖实录》洪武六年条；《明史·选举志》或《艺文志》；文集、碑刻或书目材料 互相核对；不过，译写候选以编年材料定位；实录记载反映朝廷选择，崇祯期须另以奏疏、地方及敌对政权材料交叉。 因而这里标示的是材料能够支持的行动与制度位置，而不是对动机、地方执行或普通人经验的无保留复原。

\eventanchor{ML-1373-002}{明·1373（洪武六年）十一月29日}
\paragraph{1373（洪武六年）十一月29日：明军在怀柔击败可克帖木儿} 这一节点应放回本章已经展开的权力、财政、地方社会与边地关系中理解。账本把它出现的条件概括为：前期边防、地方武装或跨区域资源调动的压力推动了该行动。 随后可追踪的变化是：它改变了后续调兵、补给或地方秩序的条件；战果和伤亡须分开核对。 可由《太祖实录》洪武六年条；《明史·兵志》及相关列传；边镇、地方志或对方材料 互相核对；不过，译写候选以编年材料定位；实录记载反映朝廷选择，崇祯期须另以奏疏、地方及敌对政权材料交叉。 因而这里标示的是材料能够支持的行动与制度位置，而不是对动机、地方执行或普通人经验的无保留复原。

\subsection{1373年前后的材料线索}

\eventanchor{ML-1373-003}{明·1373（洪武六年）}
\paragraph{1373（洪武六年）：洪武皇帝将高丽的进贡限制为每三年一次} 这一节点应放回本章已经展开的权力、财政、地方社会与边地关系中理解。账本把它出现的条件概括为：朝贡名义、边境安全与港口贸易的利益使交涉持续发生。 随后可追踪的变化是：它为后续册封、互市或海防安排留下文书与跨政权比较线索。 可由《太祖实录》洪武六年条；《明史·外国传》；《朝鲜王朝实录》、琉球《历代宝案》或海外档案 互相核对；不过，译写候选以编年材料定位；实录记载反映朝廷选择，崇祯期须另以奏疏、地方及敌对政权材料交叉。 因而这里标示的是材料能够支持的行动与制度位置，而不是对动机、地方执行或普通人经验的无保留复原。

\eventanchor{ML-1373-004}{明·1373（洪武六年）}
\paragraph{1373（洪武六年）：明朝官员制定了中国历史上第一部《宅法》} 这一节点应放回本章已经展开的权力、财政、地方社会与边地关系中理解。账本把它出现的条件概括为：继承、官僚协调或皇权运作的压力使问题进入朝廷程序。 随后可追踪的变化是：它影响后续人事与政策议程；动机和派系标签不能由单一叙事裁定。 可由《太祖实录》洪武六年条；《明史·本纪》及相关列传；诏令、奏疏或会典版本 互相核对；不过，译写候选以编年材料定位；实录记载反映朝廷选择，崇祯期须另以奏疏、地方及敌对政权材料交叉。 因而这里标示的是材料能够支持的行动与制度位置，而不是对动机、地方执行或普通人经验的无保留复原。

\eventanchor{ML-1374-001}{明·1374（洪武七年）}
\paragraph{1374（洪武七年）：朝廷围绕卫所军户的编籍、屯田与调发整理本年军政文书，作为明初军事接收的可核查节点。} 这一节点应放回本章已经展开的权力、财政、地方社会与边地关系中理解。账本把它出现的条件概括为：前期边防、地方武装或跨区域资源调动的压力推动了该行动。 随后可追踪的变化是：它改变了后续调兵、补给或地方秩序的条件；战果和伤亡须分开核对。 可由《太祖实录》洪武七年条；《明史·兵志》及相关列传；边镇、地方志或对方材料 互相核对；不过，桥接条目只标示可回查的年度问题与材料链；实录有中央选择性，崇祯期尤其需要奏疏、地方、朝鲜及满文材料交叉。 因而这里标示的是材料能够支持的行动与制度位置，而不是对动机、地方执行或普通人经验的无保留复原。

\eventanchor{ML-1374-002}{明·1374（洪武七年）}
\paragraph{1374（洪武七年）：北元诸部、东北和西南的边报、招抚与防御命令持续进入本年政务，须按具体部众和地域复核。} 这一节点应放回本章已经展开的权力、财政、地方社会与边地关系中理解。账本把它出现的条件概括为：朝贡名义、边境安全与港口贸易的利益使交涉持续发生。 随后可追踪的变化是：它为后续册封、互市或海防安排留下文书与跨政权比较线索。 可由《太祖实录》洪武七年条；《明史·外国传》；《朝鲜王朝实录》、琉球《历代宝案》或海外档案 互相核对；不过，桥接条目只标示可回查的年度问题与材料链；实录有中央选择性，崇祯期尤其需要奏疏、地方、朝鲜及满文材料交叉。 因而这里标示的是材料能够支持的行动与制度位置，而不是对动机、地方执行或普通人经验的无保留复原。

\eventanchor{ML-1374-003}{明·1374（洪武七年）}
\paragraph{1374（洪武七年）：沿海朝贡、海禁和港口管理的文书构成本年海上联系线索，禁令不能直接推作私人航行停止。} 这一节点应放回本章已经展开的权力、财政、地方社会与边地关系中理解。账本把它出现的条件概括为：朝贡名义、边境安全与港口贸易的利益使交涉持续发生。 随后可追踪的变化是：它为后续册封、互市或海防安排留下文书与跨政权比较线索。 可由《太祖实录》洪武七年条；《明史·外国传》；《朝鲜王朝实录》、琉球《历代宝案》或海外档案 互相核对；不过，桥接条目只标示可回查的年度问题与材料链；实录有中央选择性，崇祯期尤其需要奏疏、地方、朝鲜及满文材料交叉。 因而这里标示的是材料能够支持的行动与制度位置，而不是对动机、地方执行或普通人经验的无保留复原。

\eventanchor{EM-1376-PROVINCIAL-COMMISSIONS}{明·1376（洪武九年）六月}
\paragraph{1376（洪武九年）六月：废行中书省，分置承宣布政使司、都指挥使司、提刑按察使司} 这一节点应放回本章已经展开的权力、财政、地方社会与边地关系中理解。账本把它出现的条件概括为：太祖防范行省集军政财权于一体，并需将新征服地区纳入可分察的行政链条 随后可追踪的变化是：地方行政、军事和监察被分别归属三司；协调成本和中央派遣机制随之增加 可由《太祖实录》洪武九年六月条；《明史·职官志》；《明会典·布政司》 互相核对；不过，制度发布可靠；三司分置不等于地方权力已均衡运作 因而这里标示的是材料能够支持的行动与制度位置，而不是对动机、地方执行或普通人经验的无保留复原。

\eventanchor{EM-1380-HU-WEIYONG}{明·1380（洪武十三年）正月}
\paragraph{1380（洪武十三年）正月：胡惟庸案后中书省被废，皇帝直接统辖六部} 这一节点应放回本章已经展开的权力、财政、地方社会与边地关系中理解。账本把它出现的条件概括为：太祖对宰相权力和官僚结党高度警惕；案件成为重组中枢权力的触发点 随后可追踪的变化是：六部直接对皇帝负责；政务依赖皇帝裁断与后续内阁文书协调的矛盾加深 可由《太祖实录》洪武十三年正月条；《明史·胡惟庸传》与《职官志》；《大诰》胡惟庸案条 互相核对；不过，废相与制度变动可靠；胡惟庸“通倭通元”等指控须视为案狱叙事而非既定事实 因而这里标示的是材料能够支持的行动与制度位置，而不是对动机、地方执行或普通人经验的无保留复原。

\eventanchor{EM-1381-LIJIA-HUANGCE}{明·1381（洪武十四年）}
\paragraph{1381（洪武十四年）：朝廷推进黄册与里甲编制，登记户口、田土和役属} 这一节点应放回本章已经展开的权力、财政、地方社会与边地关系中理解。账本把它出现的条件概括为：新朝须将赋役、军户和土地责任落实到可追查的户等；元末流亡和战争破坏使登记困难 随后可追踪的变化是：里甲成为赋役征派的基层框架；漏报、逃役和地方中介将持续塑造实际负担 可由《太祖实录》洪武十四年户籍条；《明会典·户口》；徽州等地明代黄册残卷 互相核对；不过，制度命令可靠；存世黄册仅反映部分地区与登记口径 因而这里标示的是材料能够支持的行动与制度位置，而不是对动机、地方执行或普通人经验的无保留复原。

\eventanchor{EM-1381-YUNNAN-CAMPAIGN}{明·1381（洪武十四年）}
\paragraph{1381（洪武十四年）：傅友德、蓝玉、沐英率军自川黔进攻云南} 这一节点应放回本章已经展开的权力、财政、地方社会与边地关系中理解。账本把它出现的条件概括为：梁王仍据云南并与北元保持联系；明廷欲控制西南交通、矿产和边疆名义 随后可追踪的变化是：大规模军事接收开始；军屯、卫所和土司安排将深刻改变云南地方秩序 可由《太祖实录》洪武十四年条；《明史·傅友德传》《沐英传》；《云南图经志书》云南府条 互相核对；不过，出兵和主帅可靠；军额、伤亡与各族群立场不可照录明方奏报 因而这里标示的是材料能够支持的行动与制度位置，而不是对动机、地方执行或普通人经验的无保留复原。

\eventanchor{EM-1382-YUNNAN-ANNEXATION}{明·1382（洪武十五年）}
\paragraph{1382（洪武十五年）：昆明陷落后梁王自尽，明廷在云南设卫所并逐步建置} 这一节点应放回本章已经展开的权力、财政、地方社会与边地关系中理解。账本把它出现的条件概括为：明军取得昆明后需维持补给并处理旧元官军、土司和山地社群 随后可追踪的变化是：云南成为军屯和土司并行的边疆区域；征调成本与地方反抗未因占城而结束 可由《太祖实录》洪武十五年条；《明史·云南土司传》；大理、昆明地区碑刻与地方志 互相核对；不过，核心结果可靠；地方降附过程和破坏程度需以地方材料分别核验 因而这里标示的是材料能够支持的行动与制度位置，而不是对动机、地方执行或普通人经验的无保留复原。

\eventanchor{EM-1383-JINYIWEI}{明·1383（洪武十六年）}
\paragraph{1383（洪武十六年）：设置锦衣卫，皇帝以亲军兼掌侦缉和诏狱} 这一节点应放回本章已经展开的权力、财政、地方社会与边地关系中理解。账本把它出现的条件概括为：废相后皇帝强化对官僚和军政信息的直接掌握；大案审讯需要绕开常规司法链 随后可追踪的变化是：亲军侦缉成为中枢权力工具之一；司法程序与政治案件的界线更趋紧张 可由《太祖实录》洪武十六年条；《明史·职官志》；《明会典·锦衣卫》 互相核对；不过，设卫可靠；后世将其职能倒投到洪武初年须谨慎 因而这里标示的是材料能够支持的行动与制度位置，而不是对动机、地方执行或普通人经验的无保留复原。

\eventanchor{EM-1385-DAGAO}{明·1385（洪武十八年）}
\paragraph{1385（洪武十八年）：《大诰》颁行，以案例和重刑语言约束官民} 这一节点应放回本章已经展开的权力、财政、地方社会与边地关系中理解。账本把它出现的条件概括为：太祖认为常律不足以威慑吏治和社会流动中的违法行为，并借案例塑造服从语言 随后可追踪的变化是：法外重典进入基层教化与审判语境；法条、诏令和地方执行之间的落差须另行追踪 可由《大诰》初编；《太祖实录》洪武十八年条；《明史·刑法志》 互相核对；不过，文本与颁行可靠；规范性惩戒话语不能直接换算为实际犯罪率 因而这里标示的是材料能够支持的行动与制度位置，而不是对动机、地方执行或普通人经验的无保留复原。

\eventanchor{EM-1385-GUO-HUAN}{明·1385（洪武十八年）}
\paragraph{1385（洪武十八年）：郭桓案牵连户部及各地征粮官吏，朝廷严惩侵没赋税者} 这一节点应放回本章已经展开的权力、财政、地方社会与边地关系中理解。账本把它出现的条件概括为：江南赋粮征解经多层胥吏和官员转运，中央难以核对定额、耗损与实收 随后可追踪的变化是：恐惧性稽核整顿财政链条但也冲击地方行政；账册真实度与官吏自保问题加剧 可由《太祖实录》洪武十八年条；《明史·刑法志》；《大诰》郭桓案条 互相核对；不过，案件与惩处可靠；被杀人数及“全国贪墨”规模高度依赖官方案狱口径 因而这里标示的是材料能够支持的行动与制度位置，而不是对动机、地方执行或普通人经验的无保留复原。

\eventanchor{EM-1387-FISHSCALE-REGISTERS}{明·1387（洪武二十年）}
\paragraph{1387（洪武二十年）：朝廷催成鱼鳞图册等田土登记以配合赋役核算} 这一节点应放回本章已经展开的权力、财政、地方社会与边地关系中理解。账本把它出现的条件概括为：黄册登记户等后，国家需以地块和赋额关联田土责任；土地转移与隐匿造成核算压力 随后可追踪的变化是：田土登记成为赋役分配的凭据；地方书手、业主和契约实践影响其可执行性 可由《太祖实录》洪武二十年田土条；《明会典·田土》；徽州鱼鳞图册与相关契约 互相核对；不过，命令与部分册籍可证；图册年代、抄补层次和亩数不等于实际占有 因而这里标示的是材料能够支持的行动与制度位置，而不是对动机、地方执行或普通人经验的无保留复原。

\eventanchor{EM-1387-NAGHACHU}{明·1387（洪武二十年）}
\paragraph{1387（洪武二十年）：冯胜等在辽东、松花江方向迫使纳哈出率部降明} 这一节点应放回本章已经展开的权力、财政、地方社会与边地关系中理解。账本把它出现的条件概括为：北元残部在辽东与松花江流域保持机动空间；明廷需稳住山海关以东通道 随后可追踪的变化是：明在辽东增强卫所布置与招抚能力；东北诸部的朝贡和军事关系仍不断变动 可由《太祖实录》洪武二十年条；《明史·冯胜传》与《北虏传》；《高丽史》辽东方面条 互相核对；不过，降附事实可靠；部众数量和“归化”程度不能由明方献俘叙事推定 因而这里标示的是材料能够支持的行动与制度位置，而不是对动机、地方执行或普通人经验的无保留复原。

\eventanchor{EM-1388-BUIR}{明·1388（洪武二十一年）}
\paragraph{1388（洪武二十一年）：蓝玉军在捕鱼儿海击溃北元主力并俘获大批人员物资} 这一节点应放回本章已经展开的权力、财政、地方社会与边地关系中理解。账本把它出现的条件概括为：明军利用北元内部不稳和长距离奔袭寻求决定性战果 随后可追踪的变化是：北元汗庭进一步分散；蓝玉军功骤升，也加深洪武朝对勋贵兵权的疑惧 可由《太祖实录》洪武二十一年条；《明史·蓝玉传》与《北虏传》；17世纪《蒙古源流》北元汗统叙事 互相核对；不过，胜利可靠；战果数字、元主脱逃经过和俘获规模应保留两方材料差异 因而这里标示的是材料能够支持的行动与制度位置，而不是对动机、地方执行或普通人经验的无保留复原。

\eventanchor{EM-1390-LAN-YU}{明·1390（洪武二十三年）}
\paragraph{1390（洪武二十三年）：蓝玉以谋反罪被诛，功臣集团再受清洗} 这一节点应放回本章已经展开的权力、财政、地方社会与边地关系中理解。账本把它出现的条件概括为：北伐功臣掌握军中网络且太祖持续防范结党；案狱成为重排军政人事的渠道 随后可追踪的变化是：开国武将集团受到重创；皇太孙继承安排更依赖文官和制度性控制 可由《太祖实录》洪武二十三年条；《明史·蓝玉传》；《大诰》相关案件 互相核对；不过，处决与案名可靠；谋反细节主要来自后修实录和官方传记，难以独立裁定 因而这里标示的是材料能够支持的行动与制度位置，而不是对动机、地方执行或普通人经验的无保留复原。

\eventanchor{EM-1392-HEIR-CHANGE}{明·1392（洪武二十五年）四月}
\paragraph{1392（洪武二十五年）四月：皇太子朱标去世，太祖立朱允炆为皇太孙} 这一节点应放回本章已经展开的权力、财政、地方社会与边地关系中理解。账本把它出现的条件概括为：朱标早逝打破原有继承顺序；太祖需在皇孙与成年藩王之间确立可被承认的名分 随后可追踪的变化是：建文君臣日后削藩拥有继承政治背景；燕王等边王的资源与疑虑同时上升 可由《太祖实录》洪武二十五年四月条；《明史·懿文太子传》与《惠帝本纪》；南京懿文太子陵考古简报与碑刻著录 互相核对；不过，继承变更可靠；对朱标能力和诸王态度的评价多为后出叙事 因而这里标示的是材料能够支持的行动与制度位置，而不是对动机、地方执行或普通人经验的无保留复原。

\eventanchor{EM-1393-BLUE-JADE-PURGE}{明·1393（洪武二十六年）}
\paragraph{1393（洪武二十六年）：蓝玉案牵连扩大，勋贵与官员大批获罪} 这一节点应放回本章已经展开的权力、财政、地方社会与边地关系中理解。账本把它出现的条件概括为：太祖以案狱消除潜在军事和官僚自主网络，并为皇太孙预作权力清场 随后可追踪的变化是：朝廷人事更加依附皇权；以高压保障继承也削弱了可制衡藩王的开国武臣 可由《太祖实录》洪武二十六年条；《明史·功臣世表》与相关列传；《大诰三编》 互相核对；不过，清洗持续及其政治影响可靠；受牵连人数和罪证须避免采纳单一官方总数 因而这里标示的是材料能够支持的行动与制度位置，而不是对动机、地方执行或普通人经验的无保留复原。

\eventanchor{EM-1397-DAMING-LAW}{明·1397（洪武三十年）}
\paragraph{1397（洪武三十年）：定颁《大明律》终本，形成以六律为纲的成文法框架} 这一节点应放回本章已经展开的权力、财政、地方社会与边地关系中理解。账本把它出现的条件概括为：洪武前期反复修律与《大诰》重典并行，中央需要稳定、可传抄的法典文本 随后可追踪的变化是：终本成为明代司法基本参照；诏令、条例和诏狱仍会改变实际司法面貌 可由《大明律》洪武三十年刊本系统；《太祖实录》洪武三十年条；《明史·刑法志》 互相核对；不过，终本年代可靠；法典条文不能证明州县审判或社会遵从一致 因而这里标示的是材料能够支持的行动与制度位置，而不是对动机、地方执行或普通人经验的无保留复原。

\subsection{1397年前后的材料线索}

\eventanchor{EM-1397-SOUTHERN-LIST}{明·1397（洪武三十年）}
\paragraph{1397（洪武三十年）：会试南北榜案后重取中式者，籍贯分布被纳入取士政治} 这一节点应放回本章已经展开的权力、财政、地方社会与边地关系中理解。账本把它出现的条件概括为：江南考生在考试中占多，北方士人及皇帝对选才地域失衡产生质疑 随后可追踪的变化是：取士过程被强化为皇权可干预的政治问题；南北地域名额与资格争论延续 可由《太祖实录》洪武三十年条；《明史·选举志》；《明会典·礼部》科举条 互相核对；不过，事件可靠；“南北地域偏私”是朝廷处理语言，不能还原为单一真实原因 因而这里标示的是材料能够支持的行动与制度位置，而不是对动机、地方执行或普通人经验的无保留复原。

\eventanchor{EM-1398-HONGWU-DEATH}{明·1398（洪武三十一年）闰五月}
\paragraph{1398（洪武三十一年）闰五月：太祖去世，皇太孙朱允炆即位为建文帝} 这一节点应放回本章已经展开的权力、财政、地方社会与边地关系中理解。账本把它出现的条件概括为：太祖已以皇太孙继承并清洗部分功臣，但成年藩王仍掌边镇卫所资源 随后可追踪的变化是：建文新政面对削藩与守成的选择；中央—燕王冲突迅速制度化 可由《太祖实录》洪武三十一年闰五月条；《明史·惠帝本纪》；孝陵碑刻与陵寝考古资料 互相核对；不过，继位日期可靠；遗诏与宫廷部署的细节受建文、永乐两朝修史影响 因而这里标示的是材料能够支持的行动与制度位置，而不是对动机、地方执行或普通人经验的无保留复原。

\eventanchor{EM-1399-JINGNAN-BEGIN}{明·1399（建文元年）七月}
\paragraph{1399（建文元年）七月：燕王朱棣以“靖难”为名起兵反对建文朝} 这一节点应放回本章已经展开的权力、财政、地方社会与边地关系中理解。账本把它出现的条件概括为：建文朝先后处理周、齐、湘等王，燕王担忧自身兵权和封国安全 随后可追踪的变化是：内战使北方卫所、运河和山东交通成为争夺对象；胜负将决定继承叙事的书写者 可由《明太宗实录》建文元年七月条；《明史·成祖本纪》与《方孝孺传》；北平地方碑刻和方志线索 互相核对；不过，起兵和主要诏令可靠；双方“奉天靖难”与“削藩”合法性措辞均具政治目的 因而这里标示的是材料能够支持的行动与制度位置，而不是对动机、地方执行或普通人经验的无保留复原。

\eventanchor{EM-1400-BAIGOU}{明·1400（建文二年）}
\paragraph{1400（建文二年）：白沟河战役后燕军突破南下通道} 这一节点应放回本章已经展开的权力、财政、地方社会与边地关系中理解。账本把它出现的条件概括为：建文军依赖李景隆等集中兵力围燕；燕军须夺取河北南下道路和补给 随后可追踪的变化是：燕军获得向山东、直隶推进的机会；中央军统帅和兵源组织的弱点暴露 可由《明太宗实录》建文二年条；《明史·成祖本纪》与《李景隆传》；北平、保定地方志战场条 互相核对；不过，战役结局大体可靠；火器作用、军数和将领责任在双方叙事中夸张不一 因而这里标示的是材料能够支持的行动与制度位置，而不是对动机、地方执行或普通人经验的无保留复原。

\eventanchor{EM-1401-DONGCHANG}{明·1401（建文三年）十二月}
\paragraph{1401（建文三年）十二月：东昌之战中燕军受挫，张玉战死} 这一节点应放回本章已经展开的权力、财政、地方社会与边地关系中理解。账本把它出现的条件概括为：燕军南下后补给线拉长，建文军在山东仍可集结守军反击 随后可追踪的变化是：燕军一度受阻并调整战略；持续内战增加沿线州县征发和交通损耗 可由《明太宗实录》建文三年十二月条；《明史·张玉传》；山东地方志东昌战事条 互相核对；不过，战役与张玉死亡可靠；胜败归因及伤亡数字不宜只用《太宗实录》 因而这里标示的是材料能够支持的行动与制度位置，而不是对动机、地方执行或普通人经验的无保留复原。

\eventanchor{EM-1402-NANJING-FALL}{明·1402（建文四年）六月}
\paragraph{1402（建文四年）六月：燕军入南京，建文帝去向不明，朱棣即位} 这一节点应放回本章已经展开的权力、财政、地方社会与边地关系中理解。账本把它出现的条件概括为：燕军控制长江北岸与运河通道，建文朝内部指挥和防御体系瓦解 随后可追踪的变化是：永乐朝开始重编前朝记录并重置人事；建文遗臣与继承合法性问题长期存在 可由《明太宗实录》建文四年六月条；《明史·成祖本纪》《惠帝本纪》；南京明故宫考古与地方传说材料 互相核对；不过，入城和即位可靠；宫城火灾、建文帝生死及出逃传说不能裁作确定事实 因而这里标示的是材料能够支持的行动与制度位置，而不是对动机、地方执行或普通人经验的无保留复原。

\eventanchor{ML-1402-006}{明·1402（建文四年）六月7日}
\paragraph{1402（建文四年）六月7日：靖难之役：朱棣大军渡过淮河} 这一节点应放回本章已经展开的权力、财政、地方社会与边地关系中理解。账本把它出现的条件概括为：财政征解、交通工程或货币支付的压力使相关措施进入政务。 随后可追踪的变化是：其法定安排不等于地方执行，后续须以地域材料追踪负担与成效。 可由《惠帝实录》建文四年条；《明会典》相关条目；地方志、黄册/契约/河工材料 互相核对；不过，译写候选以编年材料定位；实录记载反映朝廷选择，崇祯期须另以奏疏、地方及敌对政权材料交叉。 因而这里标示的是材料能够支持的行动与制度位置，而不是对动机、地方执行或普通人经验的无保留复原。

\eventanchor{ML-1402-007}{明·1402（建文四年）六月17日}
\paragraph{1402（建文四年）六月17日：靖难之役：朱棣拿下扬州} 这一节点应放回本章已经展开的权力、财政、地方社会与边地关系中理解。账本把它出现的条件概括为：前期边防、地方武装或跨区域资源调动的压力推动了该行动。 随后可追踪的变化是：它改变了后续调兵、补给或地方秩序的条件；战果和伤亡须分开核对。 可由《惠帝实录》建文四年条；《明史·兵志》及相关列传；边镇、地方志或对方材料 互相核对；不过，译写候选以编年材料定位；实录记载反映朝廷选择，崇祯期须另以奏疏、地方及敌对政权材料交叉。 因而这里标示的是材料能够支持的行动与制度位置，而不是对动机、地方执行或普通人经验的无保留复原。

\eventanchor{ML-1402-008}{明·1402（建文四年）七月1日}
\paragraph{1402（建文四年）七月1日：靖难之役：朱棣被拦在南京对面的长江边} 这一节点应放回本章已经展开的权力、财政、地方社会与边地关系中理解。账本把它出现的条件概括为：前期边防、地方武装或跨区域资源调动的压力推动了该行动。 随后可追踪的变化是：它改变了后续调兵、补给或地方秩序的条件；战果和伤亡须分开核对。 可由《惠帝实录》建文四年条；《明史·兵志》及相关列传；边镇、地方志或对方材料 互相核对；不过，译写候选以编年材料定位；实录记载反映朝廷选择，崇祯期须另以奏疏、地方及敌对政权材料交叉。 因而这里标示的是材料能够支持的行动与制度位置，而不是对动机、地方执行或普通人经验的无保留复原。

\eventanchor{ML-1402-009}{明·1402（建文四年）七月3日}
\paragraph{1402（建文四年）七月3日：靖难之役：布政司侍郎陈选投奔朱棣，起义军渡过长江} 这一节点应放回本章已经展开的权力、财政、地方社会与边地关系中理解。账本把它出现的条件概括为：前期边防、地方武装或跨区域资源调动的压力推动了该行动。 随后可追踪的变化是：它改变了后续调兵、补给或地方秩序的条件；战果和伤亡须分开核对。 可由《惠帝实录》建文四年条；《明史·兵志》及相关列传；边镇、地方志或对方材料 互相核对；不过，译写候选以编年材料定位；实录记载反映朝廷选择，崇祯期须另以奏疏、地方及敌对政权材料交叉。 因而这里标示的是材料能够支持的行动与制度位置，而不是对动机、地方执行或普通人经验的无保留复原。

\eventanchor{ML-1402-010}{明·1402（建文四年）七月13日}
\paragraph{1402（建文四年）七月13日：靖难之战：朱晖打开南京金川门，不战而让朱棣入城；建文帝失踪，家人被监禁} 这一节点应放回本章已经展开的权力、财政、地方社会与边地关系中理解。账本把它出现的条件概括为：前期边防、地方武装或跨区域资源调动的压力推动了该行动。 随后可追踪的变化是：它改变了后续调兵、补给或地方秩序的条件；战果和伤亡须分开核对。 可由《惠帝实录》建文四年条；《明史·兵志》及相关列传；边镇、地方志或对方材料 互相核对；不过，译写候选以编年材料定位；实录记载反映朝廷选择，崇祯期须另以奏疏、地方及敌对政权材料交叉。 因而这里标示的是材料能够支持的行动与制度位置，而不是对动机、地方执行或普通人经验的无保留复原。

\eventanchor{ML-1402-011}{明·1402（建文四年）七月17日}
\paragraph{1402（建文四年）七月17日：朱棣即位，是为永乐皇帝} 这一节点应放回本章已经展开的权力、财政、地方社会与边地关系中理解。账本把它出现的条件概括为：继承、官僚协调或皇权运作的压力使问题进入朝廷程序。 随后可追踪的变化是：它影响后续人事与政策议程；动机和派系标签不能由单一叙事裁定。 可由《惠帝实录》建文四年条；《明史·本纪》及相关列传；诏令、奏疏或会典版本 互相核对；不过，译写候选以编年材料定位；实录记载反映朝廷选择，崇祯期须另以奏疏、地方及敌对政权材料交叉。 因而这里标示的是材料能够支持的行动与制度位置，而不是对动机、地方执行或普通人经验的无保留复原。

\eventanchor{ML-1402-012}{明·1402（建文四年）九月}
\paragraph{1402（建文四年）九月：永乐皇帝委托编写《永乐百科全书》} 这一节点应放回本章已经展开的权力、财政、地方社会与边地关系中理解。账本把它出现的条件概括为：制度、教育或知识传播的需要推动了相关行动或文本形成。 随后可追踪的变化是：它提供制度和传播史的锚点，不能据此推断社会整体接受。 可由《惠帝实录》建文四年条；《明史·选举志》或《艺文志》；文集、碑刻或书目材料 互相核对；不过，译写候选以编年材料定位；实录记载反映朝廷选择，崇祯期须另以奏疏、地方及敌对政权材料交叉。 因而这里标示的是材料能够支持的行动与制度位置，而不是对动机、地方执行或普通人经验的无保留复原。

\eventanchor{EM-1403-BEIJING-NAMING}{明·1403（永乐元年）正月}
\paragraph{1403（永乐元年）正月：北平改称北京，永乐朝开始将其经营为北方政治中心} 这一节点应放回本章已经展开的权力、财政、地方社会与边地关系中理解。账本把它出现的条件概括为：永乐政权的军事基础在北平，且北元威胁要求皇帝更接近北方防线 随后可追踪的变化是：北京获得新的行政与象征地位；大运河、宫城和卫戍工程成为迁都的前提 可由《明太宗实录》永乐元年正月条；《明史·地理志》；北京城建考古调查资料 互相核对；不过，改名和诏令可靠；改名不等于当年完成迁都或城建 因而这里标示的是材料能够支持的行动与制度位置，而不是对动机、地方执行或普通人经验的无保留复原。

\eventanchor{EM-1403-YONGLE-DADIAN-COMMISSION}{明·1403（永乐元年）七月}
\paragraph{1403（永乐元年）七月：朝廷命解缙等编纂大型类书，后称《永乐大典》} 这一节点应放回本章已经展开的权力、财政、地方社会与边地关系中理解。账本把它出现的条件概括为：新政权需吸纳士人、整理经史与展示文化正统；内阁文臣获得组织工程的机会 随后可追踪的变化是：形成跨门类文献汇编项目；其抄写、贮藏和后世散佚另构成一条物质史 可由《明太宗实录》永乐元年条；《明史·解缙传》；《永乐大典》现存卷册题记 互相核对；不过，诏修及主编群体可靠；成书规模和保存本不能直接说明知识的社会传播范围 因而这里标示的是材料能够支持的行动与制度位置，而不是对动机、地方执行或普通人经验的无保留复原。

\eventanchor{ML-1403-001}{明·1403（永乐元年）二月}
\paragraph{1403（永乐元年）二月：永乐皇帝封北平为“北都” 北京} 这一节点应放回本章已经展开的权力、财政、地方社会与边地关系中理解。账本把它出现的条件概括为：继承、官僚协调或皇权运作的压力使问题进入朝廷程序。 随后可追踪的变化是：它影响后续人事与政策议程；动机和派系标签不能由单一叙事裁定。 可由《太宗实录》永乐元年条；《明史·本纪》及相关列传；诏令、奏疏或会典版本 互相核对；不过，译写候选以编年材料定位；实录记载反映朝廷选择，崇祯期须另以奏疏、地方及敌对政权材料交叉。 因而这里标示的是材料能够支持的行动与制度位置，而不是对动机、地方执行或普通人经验的无保留复原。

\eventanchor{ML-1403-002}{明·1403（永乐元年）四月}
\paragraph{1403（永乐元年）四月：永乐皇帝在大宁附近安置了忠诚的乌里安凯} 这一节点应放回本章已经展开的权力、财政、地方社会与边地关系中理解。账本把它出现的条件概括为：继承、官僚协调或皇权运作的压力使问题进入朝廷程序。 随后可追踪的变化是：它影响后续人事与政策议程；动机和派系标签不能由单一叙事裁定。 可由《太宗实录》永乐元年条；《明史·本纪》及相关列传；诏令、奏疏或会典版本 互相核对；不过，译写候选以编年材料定位；实录记载反映朝廷选择，崇祯期须另以奏疏、地方及敌对政权材料交叉。 因而这里标示的是材料能够支持的行动与制度位置，而不是对动机、地方执行或普通人经验的无保留复原。

\eventanchor{ML-1403-003}{明·1403（永乐元年）九月4日}
\paragraph{1403（永乐元年）九月4日：宝藏远航：下令建造200艘“远洋运输船”} 这一节点应放回本章已经展开的权力、财政、地方社会与边地关系中理解。账本把它出现的条件概括为：继承、官僚协调或皇权运作的压力使问题进入朝廷程序。 随后可追踪的变化是：它影响后续人事与政策议程；动机和派系标签不能由单一叙事裁定。 可由《太宗实录》永乐元年条；《明史·本纪》及相关列传；诏令、奏疏或会典版本 互相核对；不过，译写候选以编年材料定位；实录记载反映朝廷选择，崇祯期须另以奏疏、地方及敌对政权材料交叉。 因而这里标示的是材料能够支持的行动与制度位置，而不是对动机、地方执行或普通人经验的无保留复原。

\eventanchor{ML-1403-004}{明·1403（永乐元年）十二月}
\paragraph{1403（永乐元年）十二月：永乐皇帝设立建州卫} 这一节点应放回本章已经展开的权力、财政、地方社会与边地关系中理解。账本把它出现的条件概括为：继承、官僚协调或皇权运作的压力使问题进入朝廷程序。 随后可追踪的变化是：它影响后续人事与政策议程；动机和派系标签不能由单一叙事裁定。 可由《太宗实录》永乐元年条；《明史·本纪》及相关列传；诏令、奏疏或会典版本 互相核对；不过，译写候选以编年材料定位；实录记载反映朝廷选择，崇祯期须另以奏疏、地方及敌对政权材料交叉。 因而这里标示的是材料能够支持的行动与制度位置，而不是对动机、地方执行或普通人经验的无保留复原。

\eventanchor{ML-1403-005}{明·1403（永乐元年）}
\paragraph{1403（永乐元年）：日本对明朝中国的使团：足利义光派遣使者前往明朝，宣称自己是“您的臣民，日本国王”，并获得贸易特权} 这一节点应放回本章已经展开的权力、财政、地方社会与边地关系中理解。账本把它出现的条件概括为：朝贡名义、边境安全与港口贸易的利益使交涉持续发生。 随后可追踪的变化是：它为后续册封、互市或海防安排留下文书与跨政权比较线索。 可由《太宗实录》永乐元年条；《明史·外国传》；《朝鲜王朝实录》、琉球《历代宝案》或海外档案 互相核对；不过，译写候选以编年材料定位；实录记载反映朝廷选择，崇祯期须另以奏疏、地方及敌对政权材料交叉。 因而这里标示的是材料能够支持的行动与制度位置，而不是对动机、地方执行或普通人经验的无保留复原。

\subsection{1404年前后的材料线索}

\eventanchor{EM-1404-CROWN-PRINCE}{明·1404（永乐二年）}
\paragraph{1404（永乐二年）：朱高炽被立为皇太子，继承竞争暂获制度性裁定} 这一节点应放回本章已经展开的权力、财政、地方社会与边地关系中理解。账本把它出现的条件概括为：朱高炽、朱高煦各有不同的官僚和军功支持；永乐须固定继承秩序 随后可追踪的变化是：太子在南京监国和文官网络中积累经验；兄弟竞争仍在永乐晚年反复出现 可由《明太宗实录》永乐二年条；《明史·仁宗本纪》与《汉王高煦传》；杨士奇《东里文集》东宫奏议 互相核对；不过，册立可靠；太子与汉王的个人动机不可由后出党争叙事简单化 因而这里标示的是材料能够支持的行动与制度位置，而不是对动机、地方执行或普通人经验的无保留复原。

\eventanchor{ML-1404-001}{明·1404（永乐二年）三月}
\paragraph{1404（永乐二年）三月：宝藏远航：下令建造50艘“海船”} 这一节点应放回本章已经展开的权力、财政、地方社会与边地关系中理解。账本把它出现的条件概括为：继承、官僚协调或皇权运作的压力使问题进入朝廷程序。 随后可追踪的变化是：它影响后续人事与政策议程；动机和派系标签不能由单一叙事裁定。 可由《太宗实录》永乐二年条；《明史·本纪》及相关列传；诏令、奏疏或会典版本 互相核对；不过，译写候选以编年材料定位；实录记载反映朝廷选择，崇祯期须另以奏疏、地方及敌对政权材料交叉。 因而这里标示的是材料能够支持的行动与制度位置，而不是对动机、地方执行或普通人经验的无保留复原。

\eventanchor{ML-1404-002}{明·1404（永乐二年）七月}
\paragraph{1404（永乐二年）七月：卡拉德尔的恩克帖木儿被明朝廷授予亲王称号} 这一节点应放回本章已经展开的权力、财政、地方社会与边地关系中理解。账本把它出现的条件概括为：继承、官僚协调或皇权运作的压力使问题进入朝廷程序。 随后可追踪的变化是：它影响后续人事与政策议程；动机和派系标签不能由单一叙事裁定。 可由《太宗实录》永乐二年条；《明史·本纪》及相关列传；诏令、奏疏或会典版本 互相核对；不过，译写候选以编年材料定位；实录记载反映朝廷选择，崇祯期须另以奏疏、地方及敌对政权材料交叉。 因而这里标示的是材料能够支持的行动与制度位置，而不是对动机、地方执行或普通人经验的无保留复原。

\eventanchor{ML-1404-003}{明·1404（永乐二年）十月}
\paragraph{1404（永乐二年）十月：陈添平抵达南京，请求明朝恢复陈朝王位} 这一节点应放回本章已经展开的权力、财政、地方社会与边地关系中理解。账本把它出现的条件概括为：继承、官僚协调或皇权运作的压力使问题进入朝廷程序。 随后可追踪的变化是：它影响后续人事与政策议程；动机和派系标签不能由单一叙事裁定。 可由《太宗实录》永乐二年条；《明史·本纪》及相关列传；诏令、奏疏或会典版本 互相核对；不过，译写候选以编年材料定位；实录记载反映朝廷选择，崇祯期须另以奏疏、地方及敌对政权材料交叉。 因而这里标示的是材料能够支持的行动与制度位置，而不是对动机、地方执行或普通人经验的无保留复原。

\eventanchor{ML-1404-004}{明·1404（永乐二年）十二月}
\paragraph{1404（永乐二年）十二月：帖木儿入侵明朝，但在途中去世} 这一节点应放回本章已经展开的权力、财政、地方社会与边地关系中理解。账本把它出现的条件概括为：继承、官僚协调或皇权运作的压力使问题进入朝廷程序。 随后可追踪的变化是：它影响后续人事与政策议程；动机和派系标签不能由单一叙事裁定。 可由《太宗实录》永乐二年条；《明史·本纪》及相关列传；诏令、奏疏或会典版本 互相核对；不过，译写候选以编年材料定位；实录记载反映朝廷选择，崇祯期须另以奏疏、地方及敌对政权材料交叉。 因而这里标示的是材料能够支持的行动与制度位置，而不是对动机、地方执行或普通人经验的无保留复原。

\eventanchor{ML-1404-005}{明·1404（永乐二年）}
\paragraph{1404（永乐二年）：恢复全国科举考试} 这一节点应放回本章已经展开的权力、财政、地方社会与边地关系中理解。账本把它出现的条件概括为：制度、教育或知识传播的需要推动了相关行动或文本形成。 随后可追踪的变化是：它提供制度和传播史的锚点，不能据此推断社会整体接受。 可由《太宗实录》永乐二年条；《明史·选举志》或《艺文志》；文集、碑刻或书目材料 互相核对；不过，译写候选以编年材料定位；实录记载反映朝廷选择，崇祯期须另以奏疏、地方及敌对政权材料交叉。 因而这里标示的是材料能够支持的行动与制度位置，而不是对动机、地方执行或普通人经验的无保留复原。

\eventanchor{ML-1404-006}{明·1404（永乐二年）}
\paragraph{1404（永乐二年）：山西万户搬迁北京} 这一节点应放回本章已经展开的权力、财政、地方社会与边地关系中理解。账本把它出现的条件概括为：继承、官僚协调或皇权运作的压力使问题进入朝廷程序。 随后可追踪的变化是：它影响后续人事与政策议程；动机和派系标签不能由单一叙事裁定。 可由《太宗实录》永乐二年条；《明史·本纪》及相关列传；诏令、奏疏或会典版本 互相核对；不过，译写候选以编年材料定位；实录记载反映朝廷选择，崇祯期须另以奏疏、地方及敌对政权材料交叉。 因而这里标示的是材料能够支持的行动与制度位置，而不是对动机、地方执行或普通人经验的无保留复原。

\eventanchor{EM-1405-ZHENGHE-FIRST-VOYAGE}{明·1405（永乐三年）六月}
\paragraph{1405（永乐三年）六月：郑和率舰队出航，开始第一次西洋航行} 这一节点应放回本章已经展开的权力、财政、地方社会与边地关系中理解。账本把它出现的条件概括为：永乐朝借海上使团重建朝贡秩序并寻求海外情报和政治承认 随后可追踪的变化是：官方舰队把东南港口、东南亚和印度洋使节网络连接起来；航行的财政和地方劳役成本需要另证 可由《明太宗实录》永乐三年条；《明史·郑和传》；马欢《瀛涯胜览》相关航程 互相核对；不过，出航可靠；船数、最大船型和贸易规模须区分官修记载、后出传说与物证 因而这里标示的是材料能够支持的行动与制度位置，而不是对动机、地方执行或普通人经验的无保留复原。

\eventanchor{ML-1405-001}{明·1405（永乐三年）七月11日}
\paragraph{1405（永乐三年）七月11日：宝船航行：郑和率领27800人乘坐255艘船从南京出发，其中62艘为宝船，“载有圣旨给西洋各国，并根据身份等级，向其国王赠送金锦、花纹丝绸、彩纱”。} 这一节点应放回本章已经展开的权力、财政、地方社会与边地关系中理解。账本把它出现的条件概括为：继承、官僚协调或皇权运作的压力使问题进入朝廷程序。 随后可追踪的变化是：它影响后续人事与政策议程；动机和派系标签不能由单一叙事裁定。 可由《太宗实录》永乐三年条；《明史·本纪》及相关列传；诏令、奏疏或会典版本 互相核对；不过，译写候选以编年材料定位；实录记载反映朝廷选择，崇祯期须另以奏疏、地方及敌对政权材料交叉。 因而这里标示的是材料能够支持的行动与制度位置，而不是对动机、地方执行或普通人经验的无保留复原。

\eventanchor{ML-1405-002}{明·1405（永乐三年）}
\paragraph{1405（永乐三年）：北京新宫殿建筑开工} 这一节点应放回本章已经展开的权力、财政、地方社会与边地关系中理解。账本把它出现的条件概括为：继承、官僚协调或皇权运作的压力使问题进入朝廷程序。 随后可追踪的变化是：它影响后续人事与政策议程；动机和派系标签不能由单一叙事裁定。 可由《太宗实录》永乐三年条；《明史·本纪》及相关列传；诏令、奏疏或会典版本 互相核对；不过，译写候选以编年材料定位；实录记载反映朝廷选择，崇祯期须另以奏疏、地方及敌对政权材料交叉。 因而这里标示的是材料能够支持的行动与制度位置，而不是对动机、地方执行或普通人经验的无保留复原。

\eventanchor{ML-1405-003}{明·1405（永乐三年）}
\paragraph{1405（永乐三年）：北方营建、漕运与军需征解的本年文书显示资源调动压力，报数、起运和实际到达不得混同。} 这一节点应放回本章已经展开的权力、财政、地方社会与边地关系中理解。账本把它出现的条件概括为：财政征解、交通工程或货币支付的压力使相关措施进入政务。 随后可追踪的变化是：其法定安排不等于地方执行，后续须以地域材料追踪负担与成效。 可由《太宗实录》永乐三年条；《明会典》相关条目；地方志、黄册/契约/河工材料 互相核对；不过，桥接条目只标示可回查的年度问题与材料链；实录有中央选择性，崇祯期尤其需要奏疏、地方、朝鲜及满文材料交叉。 因而这里标示的是材料能够支持的行动与制度位置，而不是对动机、地方执行或普通人经验的无保留复原。

\eventanchor{EM-1406-ANNAM-INVASION}{明·1406（永乐四年）}
\paragraph{1406（永乐四年）：明以恢复陈氏为名出兵安南，进攻胡季犛政权} 这一节点应放回本章已经展开的权力、财政、地方社会与边地关系中理解。账本把它出现的条件概括为：胡氏取代陈朝引发流亡者求援；永乐朝将宗藩名义与南方扩张结合 随后可追踪的变化是：战争迅速占领主要城市，却把明军置于持续的地方抵抗和补给压力之中 可由《明太宗实录》永乐四年条；《明史·安南传》；越南《大越史记全书》相关纪年 互相核对；不过，出兵和政治名义可靠；当地支持度与战场损失不能由明方檄文推断 因而这里标示的是材料能够支持的行动与制度位置，而不是对动机、地方执行或普通人经验的无保留复原。

\eventanchor{ML-1406-001}{明·1406（永乐四年）四月4日}
\paragraph{1406（永乐四年）四月4日：陈添平和他的明朝护卫在进入琅山时遭到伏击并被杀} 这一节点应放回本章已经展开的权力、财政、地方社会与边地关系中理解。账本把它出现的条件概括为：继承、官僚协调或皇权运作的压力使问题进入朝廷程序。 随后可追踪的变化是：它影响后续人事与政策议程；动机和派系标签不能由单一叙事裁定。 可由《太宗实录》永乐四年条；《明史·本纪》及相关列传；诏令、奏疏或会典版本 互相核对；不过，译写候选以编年材料定位；实录记载反映朝廷选择，崇祯期须另以奏疏、地方及敌对政权材料交叉。 因而这里标示的是材料能够支持的行动与制度位置，而不是对动机、地方执行或普通人经验的无保留复原。

\eventanchor{ML-1406-002}{明·1406（永乐四年）十一月19日}
\paragraph{1406（永乐四年）十一月19日：明胡战争：明军入侵大五} 这一节点应放回本章已经展开的权力、财政、地方社会与边地关系中理解。账本把它出现的条件概括为：前期边防、地方武装或跨区域资源调动的压力推动了该行动。 随后可追踪的变化是：它改变了后续调兵、补给或地方秩序的条件；战果和伤亡须分开核对。 可由《太宗实录》永乐四年条；《明史·兵志》及相关列传；边镇、地方志或对方材料 互相核对；不过，译写候选以编年材料定位；实录记载反映朝廷选择，崇祯期须另以奏疏、地方及敌对政权材料交叉。 因而这里标示的是材料能够支持的行动与制度位置，而不是对动机、地方执行或普通人经验的无保留复原。

\eventanchor{ML-1406-003}{明·1406（永乐四年）十二月13日}
\paragraph{1406（永乐四年）十二月13日：明胡战争：明军占领达邦和升龙} 这一节点应放回本章已经展开的权力、财政、地方社会与边地关系中理解。账本把它出现的条件概括为：前期边防、地方武装或跨区域资源调动的压力推动了该行动。 随后可追踪的变化是：它改变了后续调兵、补给或地方秩序的条件；战果和伤亡须分开核对。 可由《太宗实录》永乐四年条；《明史·兵志》及相关列传；边镇、地方志或对方材料 互相核对；不过，译写候选以编年材料定位；实录记载反映朝廷选择，崇祯期须另以奏疏、地方及敌对政权材料交叉。 因而这里标示的是材料能够支持的行动与制度位置，而不是对动机、地方执行或普通人经验的无保留复原。

\eventanchor{ML-1406-004}{明·1406（永乐四年）}
\paragraph{1406（永乐四年）：宝藏航行：宝藏船队访问马六甲和爪哇，然后沿马六甲海峡前往阿鲁、萨穆德拉巴赛苏丹国和兰布里，那里的人民被描述为“非常诚实和真诚”，从那里出发 3 天到达安达曼群岛，然后再用 8 天到达锡兰西海岸，那里的国王充满敌意。舰队启程前往被誉为“西洋大国”的卡利卡特} 这一节点应放回本章已经展开的权力、财政、地方社会与边地关系中理解。账本把它出现的条件概括为：朝贡名义、边境安全与港口贸易的利益使交涉持续发生。 随后可追踪的变化是：它为后续册封、互市或海防安排留下文书与跨政权比较线索。 可由《太宗实录》永乐四年条；《明史·外国传》；《朝鲜王朝实录》、琉球《历代宝案》或海外档案 互相核对；不过，译写候选以编年材料定位；实录记载反映朝廷选择，崇祯期须另以奏疏、地方及敌对政权材料交叉。 因而这里标示的是材料能够支持的行动与制度位置，而不是对动机、地方执行或普通人经验的无保留复原。

\eventanchor{EM-1407-JIAOZHI-PROVINCE}{明·1407（永乐五年）}
\paragraph{1407（永乐五年）：明废大虞并设交趾承宣布政使司等机构} 这一节点应放回本章已经展开的权力、财政、地方社会与边地关系中理解。账本把它出现的条件概括为：明军俘获胡氏父子后选择直接统治而非扶立陈氏后裔 随后可追踪的变化是：官署、卫所和税役被移植到交趾；高成本占领激化地方抵抗并牵制南方军力 可由《明太宗实录》永乐五年条；《明史·地理志》与《安南传》；越南《大越史记全书》 互相核对；不过，设治诏令可靠；行政建置不能证明红河三角洲和山区已被有效控制 因而这里标示的是材料能够支持的行动与制度位置，而不是对动机、地方执行或普通人经验的无保留复原。

\eventanchor{EM-1407-ZHENGHE-SECOND-VOYAGE}{明·1407（永乐五年）}
\paragraph{1407（永乐五年）：郑和第二次航行出发，护送各国使节返航} 这一节点应放回本章已经展开的权力、财政、地方社会与边地关系中理解。账本把它出现的条件概括为：第一次航行带回多国使者，明廷以返送维持官方互惠与海上声望 随后可追踪的变化是：航行由一次远征转为可重复的使节循环；港口补给和翻译者作用更加关键 可由《明太宗实录》永乐五年条；《明史·郑和传》；费信《星槎胜览》 互相核对；不过，航行存在可靠；访问港口的先后和个别使团身份需与碑刻、海外材料互校 因而这里标示的是材料能够支持的行动与制度位置，而不是对动机、地方执行或普通人经验的无保留复原。

\eventanchor{ML-1407-001}{明·1407（永乐五年）}
\paragraph{1407（永乐五年）：宝藏航行：宝藏船队在巨港击败了陈祖义的海盗船队，并任命石进清为“统治当地原住民的大酋长”} 这一节点应放回本章已经展开的权力、财政、地方社会与边地关系中理解。账本把它出现的条件概括为：前期边防、地方武装或跨区域资源调动的压力推动了该行动。 随后可追踪的变化是：它改变了后续调兵、补给或地方秩序的条件；战果和伤亡须分开核对。 可由《太宗实录》永乐五年条；《明史·兵志》及相关列传；边镇、地方志或对方材料 互相核对；不过，译写候选以编年材料定位；实录记载反映朝廷选择，崇祯期须另以奏疏、地方及敌对政权材料交叉。 因而这里标示的是材料能够支持的行动与制度位置，而不是对动机、地方执行或普通人经验的无保留复原。

\eventanchor{ML-1407-002}{明·1407（永乐五年）三月13日}
\paragraph{1407（永乐五年）三月13日：明—大五（胡朝）战争：胡曲里对明军的反攻失败} 这一节点应放回本章已经展开的权力、财政、地方社会与边地关系中理解。账本把它出现的条件概括为：前期边防、地方武装或跨区域资源调动的压力推动了该行动。 随后可追踪的变化是：它改变了后续调兵、补给或地方秩序的条件；战果和伤亡须分开核对。 可由《太宗实录》永乐五年条；《明史·兵志》及相关列传；边镇、地方志或对方材料 互相核对；不过，译写候选以编年材料定位；实录记载反映朝廷选择，崇祯期须另以奏疏、地方及敌对政权材料交叉。 因而这里标示的是材料能够支持的行动与制度位置，而不是对动机、地方执行或普通人经验的无保留复原。

\subsection{1407年前后的材料线索}

\eventanchor{ML-1407-003}{明·1407（永乐五年）四月}
\paragraph{1407（永乐五年）四月：五世噶玛巴喇嘛德新谢巴抵达南京举行宗教仪式} 这一节点应放回本章已经展开的权力、财政、地方社会与边地关系中理解。账本把它出现的条件概括为：继承、官僚协调或皇权运作的压力使问题进入朝廷程序。 随后可追踪的变化是：它影响后续人事与政策议程；动机和派系标签不能由单一叙事裁定。 可由《太宗实录》永乐五年条；《明史·本纪》及相关列传；诏令、奏疏或会典版本 互相核对；不过，译写候选以编年材料定位；实录记载反映朝廷选择，崇祯期须另以奏疏、地方及敌对政权材料交叉。 因而这里标示的是材料能够支持的行动与制度位置，而不是对动机、地方执行或普通人经验的无保留复原。

\eventanchor{ML-1407-004}{明·1407（永乐五年）六月16日}
\paragraph{1407（永乐五年）六月16日：明—大五（胡朝）战争：胡贵利父子被俘，送往南京} 这一节点应放回本章已经展开的权力、财政、地方社会与边地关系中理解。账本把它出现的条件概括为：前期边防、地方武装或跨区域资源调动的压力推动了该行动。 随后可追踪的变化是：它改变了后续调兵、补给或地方秩序的条件；战果和伤亡须分开核对。 可由《太宗实录》永乐五年条；《明史·兵志》及相关列传；边镇、地方志或对方材料 互相核对；不过，译写候选以编年材料定位；实录记载反映朝廷选择，崇祯期须另以奏疏、地方及敌对政权材料交叉。 因而这里标示的是材料能够支持的行动与制度位置，而不是对动机、地方执行或普通人经验的无保留复原。

\eventanchor{ML-1407-005}{明·1407（永乐五年）七月5日}
\paragraph{1407（永乐五年）七月5日：中国第四次统治越南：永乐皇帝宣布交趾正式并入明朝} 这一节点应放回本章已经展开的权力、财政、地方社会与边地关系中理解。账本把它出现的条件概括为：朝贡名义、边境安全与港口贸易的利益使交涉持续发生。 随后可追踪的变化是：它为后续册封、互市或海防安排留下文书与跨政权比较线索。 可由《太宗实录》永乐五年条；《明史·外国传》；《朝鲜王朝实录》、琉球《历代宝案》或海外档案 互相核对；不过，译写候选以编年材料定位；实录记载反映朝廷选择，崇祯期须另以奏疏、地方及敌对政权材料交叉。 因而这里标示的是材料能够支持的行动与制度位置，而不是对动机、地方执行或普通人经验的无保留复原。

\eventanchor{ML-1407-006}{明·1407（永乐五年）十月2日}
\paragraph{1407（永乐五年）十月2日：宝藏之旅：中国宝藏船队返回南京} 这一节点应放回本章已经展开的权力、财政、地方社会与边地关系中理解。账本把它出现的条件概括为：继承、官僚协调或皇权运作的压力使问题进入朝廷程序。 随后可追踪的变化是：它影响后续人事与政策议程；动机和派系标签不能由单一叙事裁定。 可由《太宗实录》永乐五年条；《明史·本纪》及相关列传；诏令、奏疏或会典版本 互相核对；不过，译写候选以编年材料定位；实录记载反映朝廷选择，崇祯期须另以奏疏、地方及敌对政权材料交叉。 因而这里标示的是材料能够支持的行动与制度位置，而不是对动机、地方执行或普通人经验的无保留复原。

\eventanchor{ML-1407-007}{明·1407（永乐五年）十月5日}
\paragraph{1407（永乐五年）十月5日：宝藏远航：王浩奉命改装249艘“海上运输船”，为“驻西洋各国使馆做准备”} 这一节点应放回本章已经展开的权力、财政、地方社会与边地关系中理解。账本把它出现的条件概括为：继承、官僚协调或皇权运作的压力使问题进入朝廷程序。 随后可追踪的变化是：它影响后续人事与政策议程；动机和派系标签不能由单一叙事裁定。 可由《太宗实录》永乐五年条；《明史·本纪》及相关列传；诏令、奏疏或会典版本 互相核对；不过，译写候选以编年材料定位；实录记载反映朝廷选择，崇祯期须另以奏疏、地方及敌对政权材料交叉。 因而这里标示的是材料能够支持的行动与制度位置，而不是对动机、地方执行或普通人经验的无保留复原。

\eventanchor{ML-1407-008}{明·1407（永乐五年）十月30日}
\paragraph{1407（永乐五年）十月30日：宝藏之旅：一位太监总督带着一封给占婆国王的御书启程} 这一节点应放回本章已经展开的权力、财政、地方社会与边地关系中理解。账本把它出现的条件概括为：朝贡名义、边境安全与港口贸易的利益使交涉持续发生。 随后可追踪的变化是：它为后续册封、互市或海防安排留下文书与跨政权比较线索。 可由《太宗实录》永乐五年条；《明史·外国传》；《朝鲜王朝实录》、琉球《历代宝案》或海外档案 互相核对；不过，译写候选以编年材料定位；实录记载反映朝廷选择，崇祯期须另以奏疏、地方及敌对政权材料交叉。 因而这里标示的是材料能够支持的行动与制度位置，而不是对动机、地方执行或普通人经验的无保留复原。

\eventanchor{ML-1407-009}{明·1407（永乐五年）}
\paragraph{1407（永乐五年）：郑和下西洋：郑和船队由249艘船组成，航线与第一次下西洋相似，增加了嘉义勒、阿博巴丹、甘巴里、奎隆、科钦等地} 这一节点应放回本章已经展开的权力、财政、地方社会与边地关系中理解。账本把它出现的条件概括为：继承、官僚协调或皇权运作的压力使问题进入朝廷程序。 随后可追踪的变化是：它影响后续人事与政策议程；动机和派系标签不能由单一叙事裁定。 可由《太宗实录》永乐五年条；《明史·本纪》及相关列传；诏令、奏疏或会典版本 互相核对；不过，译写候选以编年材料定位；实录记载反映朝廷选择，崇祯期须另以奏疏、地方及敌对政权材料交叉。 因而这里标示的是材料能够支持的行动与制度位置，而不是对动机、地方执行或普通人经验的无保留复原。

\eventanchor{ML-1407-010}{明·1407（永乐五年）十二月}
\paragraph{1407（永乐五年）十二月：《永乐百科全书》完成} 这一节点应放回本章已经展开的权力、财政、地方社会与边地关系中理解。账本把它出现的条件概括为：制度、教育或知识传播的需要推动了相关行动或文本形成。 随后可追踪的变化是：它提供制度和传播史的锚点，不能据此推断社会整体接受。 可由《太宗实录》永乐五年条；《明史·选举志》或《艺文志》；文集、碑刻或书目材料 互相核对；不过，译写候选以编年材料定位；实录记载反映朝廷选择，崇祯期须另以奏疏、地方及敌对政权材料交叉。 因而这里标示的是材料能够支持的行动与制度位置，而不是对动机、地方执行或普通人经验的无保留复原。

\eventanchor{ML-1407-011}{明·1407（永乐五年）}
\paragraph{1407（永乐五年）：明朝大炮中添加了铁木填料，提高了其效能。} 这一节点应放回本章已经展开的权力、财政、地方社会与边地关系中理解。账本把它出现的条件概括为：制度、教育或知识传播的需要推动了相关行动或文本形成。 随后可追踪的变化是：它提供制度和传播史的锚点，不能据此推断社会整体接受。 可由《太宗实录》永乐五年条；《明史·选举志》或《艺文志》；文集、碑刻或书目材料 互相核对；不过，译写候选以编年材料定位；实录记载反映朝廷选择，崇祯期须另以奏疏、地方及敌对政权材料交叉。 因而这里标示的是材料能够支持的行动与制度位置，而不是对动机、地方执行或普通人经验的无保留复原。

\eventanchor{EM-1408-YONGLE-DADIAN-COMPLETED}{明·1408（永乐六年）}
\paragraph{1408（永乐六年）：《永乐大典》进呈，类书编纂工程完成一阶段} 这一节点应放回本章已经展开的权力、财政、地方社会与边地关系中理解。账本把它出现的条件概括为：朝廷投入大批儒臣与抄手完成分类汇编，以文化工程巩固永乐正统 随后可追踪的变化是：文献汇编成为内府收藏与后世辑佚的重要媒介；政治整肃中的编者命运也显示工程依赖宫廷权力 可由《明太宗实录》永乐六年条；《明史·解缙传》；现存《永乐大典》卷端题记 互相核对；不过，进呈可靠；“收尽天下书”是工程宣称，不能以之推定文献无遗漏 因而这里标示的是材料能够支持的行动与制度位置，而不是对动机、地方执行或普通人经验的无保留复原。

\eventanchor{ML-1408-001}{明·1408（永乐六年）二月14日}
\paragraph{1408（永乐六年）二月14日：宝藏远航：南京工部下令建造48艘宝船} 这一节点应放回本章已经展开的权力、财政、地方社会与边地关系中理解。账本把它出现的条件概括为：继承、官僚协调或皇权运作的压力使问题进入朝廷程序。 随后可追踪的变化是：它影响后续人事与政策议程；动机和派系标签不能由单一叙事裁定。 可由《太宗实录》永乐六年条；《明史·本纪》及相关列传；诏令、奏疏或会典版本 互相核对；不过，译写候选以编年材料定位；实录记载反映朝廷选择，崇祯期须另以奏疏、地方及敌对政权材料交叉。 因而这里标示的是材料能够支持的行动与制度位置，而不是对动机、地方执行或普通人经验的无保留复原。

\eventanchor{ML-1408-002}{明·1408（永乐六年）七月5日}
\paragraph{1408（永乐六年）七月5日：中国第四次统治越南：明军夺取1360万吨大米；越南有235,900头牛、牲畜和大量物资} 这一节点应放回本章已经展开的权力、财政、地方社会与边地关系中理解。账本把它出现的条件概括为：朝贡名义、边境安全与港口贸易的利益使交涉持续发生。 随后可追踪的变化是：它为后续册封、互市或海防安排留下文书与跨政权比较线索。 可由《太宗实录》永乐六年条；《明史·外国传》；《朝鲜王朝实录》、琉球《历代宝案》或海外档案 互相核对；不过，译写候选以编年材料定位；实录记载反映朝廷选择，崇祯期须另以奏疏、地方及敌对政权材料交叉。 因而这里标示的是材料能够支持的行动与制度位置，而不是对动机、地方执行或普通人经验的无保留复原。

\eventanchor{ML-1408-003}{明·1408（永乐六年）九月}
\paragraph{1408（永乐六年）九月：中国第四次统治越南：交趾的陈毅叛军} 这一节点应放回本章已经展开的权力、财政、地方社会与边地关系中理解。账本把它出现的条件概括为：朝贡名义、边境安全与港口贸易的利益使交涉持续发生。 随后可追踪的变化是：它为后续册封、互市或海防安排留下文书与跨政权比较线索。 可由《太宗实录》永乐六年条；《明史·外国传》；《朝鲜王朝实录》、琉球《历代宝案》或海外档案 互相核对；不过，译写候选以编年材料定位；实录记载反映朝廷选择，崇祯期须另以奏疏、地方及敌对政权材料交叉。 因而这里标示的是材料能够支持的行动与制度位置，而不是对动机、地方执行或普通人经验的无保留复原。
